\documentclass[11pt,letterpaper]{article}
\usepackage[utf8]{inputenc}
\usepackage[T1]{fontenc}
\usepackage{geometry}
\geometry{margin=1in}
\usepackage{hyperref}
\usepackage{setspace}
\usepackage{siunitx}
\sisetup{separate-uncertainty=true, range-phrase=--}
\onehalfspacing

\begin{document}

\begin{flushleft}
\textbf{Myke Vital Sao Temgoua} \\
Department of Physics, Faculty of Science \\
University of Yaoundé I \\
P.O. Box 812, Yaoundé, Cameroon \\
\href{mailto:myke-vital.sao@facsciences-uy1.cm}{myke-vital.sao@facsciences-uy1.cm} \\[2em]
July 5, 2026
\end{flushleft}

\vspace{1em}

\begin{flushleft}
Prof.\ Kenneth M.\ Merz, Jr., Editor-in-Chief \\
\textit{Journal of Chemical Information and Modeling} \\
American Chemical Society
\end{flushleft}

\vspace{1em}

\noindent Dear Prof.\ Merz,

We submit our revised manuscript \textbf{``Computational Discovery of Polypharmacological Antimalarial Candidates from African Natural Product Chemical Space''} for consideration as an Article in \textit{JCIM} (\textbf{Manuscript ID: ci-2026-012015}). We have comprehensively addressed the editorial comments provided by Associate Editor Prof.\ Duc Nguyen and strengthened the narrative to foreground quantitative validation over descriptive hit lists.

\textbf{Central finding: Scaffold-hopping framework validated by MMV retrodiction.} Rather than presenting a static list of \num{20} candidates, this work establishes a generative scaffold-hopping framework capable of retrospectively predicting (retrodicting) known antimalarial ground-truth hits from the MMV Malaria Box before application to the unexplored African chemical space. The MMV Malaria Box benchmark serves as our central experimental validation. Our pipeline successfully recovers MMV positive controls with strong enrichment metrics (ROC-AUC curves, early enrichment factors EF\qty{1}{\percent} and EF\qty{5}{\percent}, and BEDROC scores). Specifically, a \qty{69.8}{\percent} hit rate in the top tiers demonstrates that our framework correctly ranks known antimalarials. Because physical \textit{in vitro} testing of our novel candidates is beyond the scope of this initial computational study, this high early enrichment (EF\qty{5}{\percent}), combined with a synthetic accessibility filter (SYBA $> \num{0}$), serves as a robust orthogonal experimental proxy that supports our hit prioritization.

\textbf{Mathematical proof of scaffold novelty and safety.} To prove that the VAE-generated molecules access chemical space unreachable by classical medicinal chemistry modifications of known drugs, we performed a scaffold leap analysis. An ECFP4 nearest-neighbor search against \num{396} seed African NPs demonstrates that \qty{92.6}{\percent} of VAE-generated molecules have no neighbour above Tanimoto \num{0.4} in the seed library (mean max similarity: \num{0.206}), indicating that the generative model accesses genuinely novel chemotypes rather than mere interpolations. Scaffold-only Tanimoto analysis reveals a \num{1.84}$\times$ increase in scaffold similarity compared to whole-molecule similarity (\num{0.379} vs.\ \num{0.206}), indicating that ring systems are preserved while peripheral substituents diverge. The \num{5} control compounds from the MMV Box with experimental IC$_{50}$ and cytotoxicity data were used to calibrate our computational ADMET safety profiles.

\textbf{Scientific contribution.} This work addresses the critical bottleneck in antimalarial discovery through a three-component pipeline:
\begin{itemize}
    \item \textbf{Scaffold-aware VAE sampling:} A convolutional variational autoencoder (\num{64}-dimensional latent space) whose scaffold-preserving embeddings achieve 1.84$\times$ higher Tanimoto similarity than whole-molecule equivalents, enabling targeted coverage of chemical space while retaining synthetic accessibility.
    \item \textbf{Dual-filter consensus docking:} A robust protocol combining DiffDock and AutoDock Vina 1.2.7 against four validated \textit{P.\ falciparum} targets (PfDHFR, PfCRT, PfATP4, PfClpP), validated by DEKOIS~2.0 enrichment (EF\qty{1}{\percent} $\geq \num{2.0}$ for three targets).
    \item \textbf{Multi-Parameter Optimization (MPO):} A five-component framework integrating binding affinity, geometric confidence, drug-likeness, ADMET safety, and Lipinski compliance to identify $>\num{70}$ high-confidence polypharmacological leads.
\end{itemize}
By reducing screening requirements from $>\num{1000}$ CPU-hours per target to $<\num{10}$ CPU-hours, this framework democratizes large-scale virtual screening for research groups in malaria-endemic regions.

\textbf{Unified antimalarial design engine.} This paper constitutes the first tier of a three-tier African antimalarial design engine. The second tier (Paper~2, companion manuscript) performs resistance-aware MD/MC validation of the top-20 candidates, introducing the Resistance Resilience Score (RRS) and the African Chemical Space Index (ACSI) to classify binding stability under African-prevalent resistance mutations. The third tier (a forthcoming study on quantum-topological representations) resolves the scaffold paradox revealed by our generative expansion through the Topological Fingerprint, Tensor Network Embedding, and Quantum Kernel Score. Together, the three studies constitute the most comprehensive AI-assisted antimalarial discovery framework targeting African malaria strains published to date, grounded in African natural product chemical heritage.

\textbf{Validation and reproducibility.} Beyond the MMV retrodiction benchmark, the docking protocol was validated by redocking co-crystallized ligands (RMSD $<$ \qty{2.0}{\angstrom} for all alignable poses) and by external enrichment against DEKOIS~2.0. All validation protocols, including the MMV benchmark suite, Z-score enrichment statistics, and scaffold diversity analysis, are fully documented with open-source code. We provide our 19,913-compound generative hybrid library (Zenodo DOI: 10.5281/zenodo.19608875) as an open-access community resource to accelerate African malaria research programs.

\textbf{Scope fit.} The manuscript introduces methodological contributions in consensus docking validation, VAE-based chemical space clustering, and polypharmacological MPO scoring, all within the cheminformatics and molecular modeling scope of \textit{JCIM}. The work is not a single-target docking application; it presents a reusable framework with documented hyperparameters, public data (Zenodo DOI: 10.5281/zenodo.19608875), and open-source code.

\textbf{Suggested reviewers.}
\begin{enumerate}
    \item Dr.\ Maria Sorokina, University of Jena, Germany (\href{mailto:maria.sorokina@uni-jena.de}{maria.sorokina@uni-jena.de}) — natural product cheminformatics, COCONUT database
    \item Dr.\ Ntie-Kang Fidele, University of Buea, Cameroon (\href{mailto:fidele.ntie-kang@ubuea.cm}{fidele.ntie-kang@ubuea.cm}) — African natural product databases, virtual screening
    \item Dr.\ Gisbert Schneider, ETH Zurich, Switzerland (\href{mailto:gisbert.schneider@bsse.ethz.ch}{gisbert.schneider@bsse.ethz.ch}) — generative chemistry, scaffold analysis
    \item Dr.\ José Jiménez-Luna, Microsoft Research Cambridge, United Kingdom (\href{mailto:jjimenezluna@microsoft.com}{jjimenezluna@microsoft.com}) — deep learning for molecular modeling
    \item Dr.\ Fran\c{c}ois D\u{e}ret, Medicines for Malaria Venture, Switzerland (\href{mailto:deretf@mmv.org}{deretf@mmv.org}) — MMV Malaria Box, antimalarial screening
\end{enumerate}

\textbf{Declarations.} This manuscript is original, not under consideration elsewhere, and all authors have approved the submission. No conflicts of interest exist. Data and code are publicly available.

\vspace{1.5em}

\noindent Sincerely,

\vspace{2em}

\noindent \textbf{Myke Vital Sao Temgoua} \\
(on behalf of all co-authors)

\vspace{2em}


\end{document}
